
% Packages nécessaires
\usepackage[utf8]{inputenc}
\usepackage[T1]{fontenc}
\usepackage[french]{babel}
\usepackage{titlesec}
\usepackage{tocloft}
\usepackage{graphicx}
\usepackage{float} % pour l'option de placement [H]
\usepackage{wrapfig}   % Pour l'enroulement du texte autour des images
\usepackage{array}

% Couleurs et tableaux
\usepackage[table]{xcolor}  % Pour les couleurs de fond des lignes et en-têtes
\definecolor{headerblue}{RGB}{26, 54, 93}   % Bleu marine corporate
\definecolor{rowlight}{RGB}{245, 247, 250}   % Gris très clair pour l'alternance

% Hyperliens stylisés
\usepackage[colorlinks=true, linkcolor=headerblue, citecolor=headerblue, urlcolor=headerblue]{hyperref}

% Mathématiques et encadrés
\usepackage{amsmath, amssymb}
\usepackage{amsthm}
\usepackage{mdframed}  % Package pour les cadres

% Configuration de la page et marges
\usepackage[a4paper, left=2.5cm, right=2.5cm, top=2.5cm, bottom=2.5cm]{geometry}
\usepackage{booktabs}       % Pour des bordures nettes et professionnelles
\usepackage{tabularx}       % Pour adapter automatiquement la largeur du tableau

% En-têtes et pieds de page personnalisés (fancyhdr)
\usepackage{fancyhdr}
\setlength{\headheight}{15pt}
\pagestyle{fancy}
\fancyhf{} % Réinitialise les en-têtes et pieds de page

% En-tête : nom du chapitre à gauche et ligne fine
\fancyhead[L]{\small\itshape\nouppercase{\leftmark}}
\renewcommand{\headrulewidth}{0.4pt}

% Pied de page : ligne horizontale fine et numéro de page encadré
\renewcommand{\footrulewidth}{0.4pt}
\fancyfoot[C]{\setlength{\fboxsep}{3pt}\fbox{\footnotesize\bfseries\thepage}}

% Style pour les premières pages de chapitres (style 'plain')
\fancypagestyle{plain}{
    \fancyhf{}
    \renewcommand{\headrulewidth}{0pt}
    \renewcommand{\footrulewidth}{0.4pt}
    \fancyfoot[C]{\setlength{\fboxsep}{3pt}\fbox{\footnotesize\bfseries\thepage}}
}